\section{Introduction}
\label{sec:intro}

The introduction first establishes the role of graded phonetic similarity in
written Mandarin pronunciation learning (Section~\ref{sec:intro-background}).
It then formulates the research question (Section~\ref{sec:intro-question})
and outlines the data, supervision objectives, and representation analyses
used to address it (Section~\ref{sec:intro-approach}).

\subsection{Background}
\label{sec:intro-background}

Phonetics examines the production, acoustic properties, and perception of speech
sounds \cite{ref_phonetic_sciences}. Different speech sounds may exhibit varying
degrees of phonetic similarity, with some pronunciations being closer to each
other than others. Experimental studies demonstrate that human listeners
distinguish such differences when comparing spoken forms and individual phonemes
\cite{ref_word_similarity,ref_phoneme_similarity}. Graded phonetic similarity
therefore describes an aspect of phonetic structure that extends beyond the
identification of individual sounds. Modeling phonetic information requires
investigating not only whether a system correctly identifies pronunciations but
also whether its learned representations reflect the relationships between
different pronunciations.

We consider the learning of phonetic representations from written Chinese
character sequences, drawing inspiration from how readers associate written
characters with their pronunciations in linguistic context. Phonological
information is activated during Chinese character recognition, even when readers
perform tasks involving semantic judgments
\cite{ref_early_phonology,ref_chinese_reading}. Mandarin provides a convenient
setting for character-level pronunciation modeling, as each written character
typically corresponds to a single spoken syllable in a given context. Syllables
can be further decomposed into three phonological components: an initial, a
final, and a lexical tone \cite{ref_mandarin_phonology}. Different syllables may
share one or more phonological components, whereas distinct characters may have
identical pronunciations, resulting in substantial \textit{homophony}.
Furthermore, some characters have multiple pronunciations whose selection
depends on the surrounding linguistic context. Learning pronunciation from
written characters therefore requires a model to associate character identities
with their pronunciations while accounting for contextual information.

Computational models of Chinese reading have demonstrated that neural networks
can acquire phonological information from written characters and reproduce
patterns observed in human reading \cite{ref_yang2009}. The relative emphasis on
pronunciation and meaning during training has also been shown to affect reading
performance \cite{ref_kuo2025}. An alternative source of phonetic information is
recorded speech, from which neural encoders can extract continuous
representations that capture pronunciation differences \cite{ref_bartelds2022}.
Such acoustic information has also been incorporated into Mandarin pronunciation
prediction models \cite{ref_g2pa2024}. However, the influence of categorical
pronunciation prediction and continuous acoustic supervision, applied separately
or jointly, on graded phonetic similarity in contextualized character
representations remains insufficiently understood.

\subsection{Research Question}
\label{sec:intro-question}

This study addresses the following research question: \emph{How do different forms of phonetic supervision influence the
extent to which contextualized neural representations learned from Mandarin
character sequences reflect graded phonetic similarity?}


\subsection{Experimental Approach}
\label{sec:intro-approach}

\subsubsection{Data and Contextual Character Encoding}
\label{sec:intro-data-encoding}

We investigate phonetic representation learning using written Mandarin character
sequences and their corresponding pronunciations. The Duyu/Pinyin-Hanzi corpus
provides character sequences with contextual Pinyin annotations \cite{ref_duyu}.
Recordings of isolated Mandarin syllables are obtained from the Chinese Lessons
sound collection \cite{ref_chinese_lessons}. These sources provide the written
inputs, pronunciation labels, and acoustic reference signals used in the study.

Characters are represented at the individual character level and processed using
two contextual encoders: a bidirectional long short-term memory network (BiLSTM)
with trainable character embeddings and a pretrained Chinese BERT model.
Recurrent networks have previously been applied to Mandarin pronunciation
prediction, while BERT-based models provide an alternative approach to
contextual character representation
\cite{ref_g2pm2020,ref_bert2019,ref_g2pw2022}. Both encoders produce hidden
representations associated with character occurrences, allowing the effects of
phonetic supervision to be examined within different neural architectures.

\subsubsection{Phonetic Supervision Objectives}
\label{sec:intro-supervision}

We use Hanyu Pinyin (hereafter Pinyin), a standardized Romanization system for
Mandarin Chinese, to represent target pronunciations \cite{ref_pinyin}.
Categorical supervision is considered at two levels. The first predicts complete
Pinyin syllables as distinct classes. The second predicts initial, final, and
lexical tone separately, allowing syllables to share phonological components.
Both approaches have precedents in neural pronunciation modeling
\cite{ref_g2pm2020,ref_li_scope2022}.

Continuous acoustic supervision provides an alternative to categorical
prediction. A pretrained speech encoder extracts fixed acoustic representations
from syllable recordings \cite{ref_wav2vec2,ref_xlsr2021}. For each character
occurrence $i$, a trainable linear head projects its contextual representation
$\mathbf{h}_i$ into the acoustic target space:

\begin{equation}
\label{eq:acoustic-projection}
    \hat{\mathbf{a}}_i = W\mathbf{h}_i + \mathbf{b}.
\end{equation}

The projection is trained to match the direction of the corresponding acoustic
target $\mathbf{a}_i$ by minimizing cosine distance:

\begin{equation}
\label{eq:acoustic-loss}
    \mathcal{L}_{\mathrm{acoustic}}
    =
    \frac{1}{N}\sum_{i=1}^{N}
    \left[
    1 -
    \frac{
        \hat{\mathbf{a}}_i^\top \mathbf{a}_i
    }{
        \max(\|\hat{\mathbf{a}}_i\|_2,\epsilon)
        \max(\|\mathbf{a}_i\|_2,\epsilon)
    }
    \right],
\end{equation}

where $N$ is the number of valid character occurrences, excluding padding, and
$\epsilon=10^{-8}$ provides numerical stability. Unlike a magnitude-sensitive
regression loss, this objective encourages directional agreement with the
acoustic representations.

We compare four training objectives:

\begin{equation}
\label{eq:training-objectives}
\begin{aligned}
\mathcal{J}_{\mathrm{syllable}}
    &= \mathrm{CE}_{\mathrm{syllable}}, \\[3pt]
\mathcal{J}_{\mathrm{component}}
    &= \frac{
        \mathrm{CE}_{\mathrm{initial}}
        + \mathrm{CE}_{\mathrm{final}}
        + \mathrm{CE}_{\mathrm{tone}}
    }{3}, \\[3pt]
\mathcal{J}_{\mathrm{acoustic}}
    &= \mathcal{L}_{\mathrm{acoustic}}, \\[3pt]
\mathcal{J}_{\mathrm{combined}}
    &= \mathcal{J}_{\mathrm{component}}
       + \lambda\mathcal{L}_{\mathrm{acoustic}}.
\end{aligned}
\end{equation}

Here, each cross-entropy term is averaged over valid character occurrences, and
$\lambda$ controls the contribution of acoustic alignment in the combined
objective. The four conditions allow us to examine the effects of
complete-syllable classification, component-level prediction, continuous
acoustic supervision, and their combination on the learned representation
spaces.

\subsubsection{Assessment of Acoustic Representations}
\label{sec:intro-acoustic-assessment}

The properties of continuous supervision depend on the speech representations
used as targets. We therefore examine two pretrained speech encoders: the
original multilingual XLSR-53 model and its Mandarin-fine-tuned variant released
by Grosman. The latter has also been used to extract acoustic features from
Mandarin speech in subsequent research \cite{ref_xlsr2021,ref_grosman2021,ref_liao2026}.

We assess the extracted syllable embeddings using two complementary criteria.
The first compares their pairwise similarities with DIMSIM, an independent
measure of Mandarin phonetic distance \cite{ref_dimsim}. The second examines
whether syllables sharing additional phonological components exhibit the
expected relative similarity rankings. This analysis uses Anchor-Q, a metric
introduced in the present study and motivated by the discrimination principle of
machine ABX evaluation \cite{ref_matusevych2023}. Together, these criteria
characterize graded phonetic relationships and component-level organization in
the acoustic target spaces. The metrics do not assume that speech recognition
fine-tuning necessarily improves either property.

\subsubsection{Evaluation of Contextual Representations}
\label{sec:intro-contextual-evaluation}

The effects of supervision are evaluated using contextual hidden states
extracted from individual character occurrences. Unlike pronunciation accuracy,
which measures agreement with target labels, representation analysis examines
relationships between different pronunciations in the learned space.

We consider two pair-sampling distributions. Natural-frequency sampling reflects
the occurrence frequencies of pronunciations in the corpus, whereas
uniform-Pinyin sampling gives pronunciation types equal selection probability.
For sampled pairs, cosine similarities between contextual hidden states are
compared with the corresponding DIMSIM distances. This approach provides a
measure of how closely the representation space reflects graded Mandarin
phonetic relationships.

Anchor-Q provides a complementary evaluation of phonological component
agreement. Each comparison contains an anchor occurrence and two candidates
whose syllables differ in their agreement with the anchor on one initial, final,
or tone component. Agreement on the remaining components is held constant. A
comparison is successful when the candidate sharing the additional component has
greater cosine similarity to the anchor. The resulting score summarizes the
consistency of these relative rankings across phonological components.

DIMSIM-based analysis therefore examines correspondence with graded phonetic
distances, while Anchor-Q assesses local ranking consistency under controlled
changes in component agreement. Applying both evaluations to the four training
conditions allows us to investigate how different supervision objectives
influence the organization of contextualized Mandarin character representations.
